\documentclass[conference]{IEEEtran}
\usepackage{cite}
\usepackage{amsmath,amssymb,amsfonts}
\usepackage{algorithmic}
\usepackage{graphicx}
\usepackage{textcomp}
\usepackage{xcolor}
\usepackage{listings}
\usepackage{booktabs}
\usepackage{url}

\begin{document}

\title{PRM Framework: Process Ancestry-Based Security for HPC Environments}

\author{
\IEEEauthorblockN{}
\IEEEauthorblockA{}
}

\maketitle

\begin{abstract}
Modern High-Performance Computing (HPC) environments and cloud orchestration platforms face a critical security challenge: balancing robust protection with computational efficiency. As scientific and enterprise workloads increasingly rely on containerized environments and privileged operations, traditional security approaches prove inadequate, creating a significant gap in the security landscape. This paper presents the PRM (Policy Rule Management) Framework, a novel process ancestry-based security solution for Linux environments that provides fine-grained control over system calls. By analyzing the lineage of processes to make security decisions based on execution paths, PRM enables administrators to distinguish between identical operations initiated through different execution chains, providing unprecedented control over privileged operations without compromising system performance. Our evaluation demonstrates that PRM introduces minimal overhead (0.5-2\% for system calls) compared to traditional solutions like SELinux (5-10\%), while offering superior context awareness and container security capabilities. The framework's ability to secure containerized workloads while preserving performance makes it an ideal solution for modern HPC infrastructures where both security and computational efficiency are critical requirements.
\end{abstract}

\begin{IEEEkeywords}
security, high-performance computing, process ancestry, system calls, containers
\end{IEEEkeywords}

\section{Introduction}
Security in High-Performance Computing (HPC) environments presents unique challenges that traditional security solutions struggle to address effectively. The computational demands of scientific workloads require minimal security overhead, while the increasing adoption of containerization introduces new attack vectors that must be mitigated. Traditional security mechanisms like SELinux, AppArmor, and seccomp either introduce unacceptable performance penalties or lack the context awareness needed to make nuanced security decisions in complex HPC environments.

The "privileged operation dilemma" is particularly acute in HPC settings: how to allow legitimate privileged operations required by scientific workflows while preventing malicious exploitation. Container technologies like Docker and Singularity often require root privileges, creating significant security concerns in multi-tenant HPC clusters where isolation between users and workloads is paramount.

\subsection{Contributions}
This paper makes the following contributions:
\begin{enumerate}
\item We introduce the PRM Framework, a novel security solution that uses process ancestry analysis to make context-aware security decisions with minimal performance impact
\item We present a comprehensive architecture for system call filtering that combines efficient caching, fast-path processing, and sophisticated rule evaluation
\item We demonstrate through benchmarks that PRM introduces significantly less overhead than traditional security solutions while providing superior security capabilities
\item We provide a flexible policy framework that enables administrators to define precise security rules based on process ancestry patterns
\end{enumerate}

\subsection{Outline}
The remainder of this paper is organized as follows: Section \ref{sec:background} provides background on Linux security mechanisms and the specific challenges in HPC environments. Section \ref{sec:related} discusses related work in system call filtering and container security. Section \ref{sec:methodology} presents the methodology and architecture of the PRM Framework. Section \ref{sec:evaluation} evaluates the framework's performance and security capabilities. Section \ref{sec:discussion} discusses implications and limitations, and Section \ref{sec:conclusion} concludes with future work directions.

\section{Background}
\label{sec:background}

\subsection{Linux Security Mechanisms}
Linux security has evolved from simple discretionary access controls to sophisticated mandatory access control systems. Modern Linux security mechanisms include:

\begin{itemize}
\item \textbf{Discretionary Access Control (DAC)}: The traditional Unix permission model
\item \textbf{Mandatory Access Control (MAC)}: Systems like SELinux and AppArmor that enforce system-wide security policies
\item \textbf{Capabilities}: Fine-grained privileges that can be assigned to processes
\item \textbf{Namespaces}: Isolation mechanisms that provide process separation
\item \textbf{Seccomp}: System call filtering based on BPF rules
\end{itemize}

\subsection{HPC Security Challenges}
HPC environments present unique security challenges:

\begin{itemize}
\item \textbf{Performance Sensitivity}: Security mechanisms that introduce significant overhead can render scientific computations impractical
\item \textbf{Complex Workflows}: Scientific applications often require privileged operations that are difficult to accommodate in restrictive security models
\item \textbf{Multi-tenancy}: HPC clusters typically serve multiple users and projects, requiring strong isolation
\item \textbf{Containerization}: Increasing use of containers introduces new security concerns, particularly when privileged operations are required
\end{itemize}

\subsection{Process Ancestry and Security Context}
Process ancestry—the lineage of parent-child relationships between processes—provides valuable context for security decisions. Malicious activities often involve unusual process ancestry patterns that can be detected and blocked. Traditional security models typically lack this contextual awareness, focusing instead on static permissions or labels.

\section{Related Work}
\label{sec:related}

\subsection{System Call Filtering Approaches}
System call filtering has been explored extensively in the literature:

\begin{itemize}
\item \textbf{Systrace} \cite{provos2003improving}: One of the early systems for enforcing system call policies
\item \textbf{Seccomp-BPF} \cite{edge2015seccomp}: The standard Linux system call filtering mechanism using Berkeley Packet Filter
\item \textbf{SPEAKER} \cite{li2018speaker}: A system that uses execution path analysis for security enforcement
\item \textbf{MBOX} \cite{kim2013practical}: A lightweight sandboxing mechanism for untrusted applications
\end{itemize}

\subsection{Container Security Solutions}
Container security has received significant attention:

\begin{itemize}
\item \textbf{Docker Security} \cite{bui2015analysis}: Built-in security features including seccomp profiles and capabilities
\item \textbf{gVisor} \cite{google2018gvisor}: A container runtime that provides an additional layer of isolation
\item \textbf{Kata Containers} \cite{kata2018architecture}: A secure container runtime using lightweight VMs
\item \textbf{Singularity} \cite{kurtzer2017singularity}: A container platform designed specifically for HPC environments with security considerations
\end{itemize}

\subsection{HPC Security Research}
Security in HPC environments has been addressed by several researchers:

\begin{itemize}
\item \textbf{SUPERCLOUD} \cite{wailly2018supercloud}: A security framework for multi-cloud environments
\item \textbf{HPCG} \cite{nist2018security}: HPC security guidelines for government facilities
\item \textbf{ClusterSec} \cite{johnson2020clustersec}: A comprehensive security framework for HPC clusters
\end{itemize}